% §10.3 HyperLogLog
%
% Cross-reference resolution status:
%
% DEFINED here:
%   \label{sec:hyperloglog}                       -- this section
%   \label{def:hyperloglog}                       -- the construction
%   \label{alg:hyperloglog-update-estimate}       -- update/estimate pseudocode
%   \label{prop:hyperloglog-error}                -- standard error proposition
%
% RESOLVED (chapter 9):
%   \Cref{ch:approximate-functions}               -- chapter 9
%   \Cref{def:n-ary-channel}                      -- §9.2
%   \Cref{thm:n-ary-kronecker}                    -- §9.3
%
% RESOLVED (chapter 10, prior sections):
%   \Cref{sec:sketches-as-bernoulli}              -- §10.1
%   \Cref{sec:countmin}                           -- §10.2
%   \Cref{thm:countmin-bound}                     -- §10.2
%
% FORWARD REFS within chapter 10 (resolve in Plan 10 Task 5):
%   \Cref{sec:minhash-lsh}                        -- §10.4
%
% BIBLIOGRAPHIC (resolved in §10.6, Plan 10 Task 7):
%   flajolet-et-al-2007                           -- HyperLogLog original

\section{HyperLogLog}
\label{sec:hyperloglog}

\subsection*{Problem setup}

The streaming question of this section is complementary to the
frequency-estimation question of \Cref{sec:countmin}. Given a stream
$x_1, x_2, \ldots, x_N$ of items from a universe $\Universe$, possibly
with repetition, estimate the cardinality
\[
  n \;=\; |\{x_1, x_2, \ldots, x_N\}|
\]
of the set of distinct items seen, ignoring multiplicities. Exact
counting requires space that grows with $n$: any deterministic
algorithm that reports the exact cardinality of an arbitrary stream
must, in essence, remember which distinct items it has already
encountered, and a standard communication-complexity argument
establishes that this is $\Theta(n)$ bits in the worst case. The
HyperLogLog sketch gives up exactness and returns an estimate $\hat n$
with relative error on the order of $1/\sqrt{m}$, using $m$
small-integer registers and a single hash evaluation per stream item.
The resulting space cost, for fixed target accuracy, is
\emph{doubly logarithmic} in $n$: a few kilobytes of state suffices
to estimate cardinalities into the billions.

\subsection*{The leading-zero trick}

Let $h : \Universe \to [0, 1]$ be a uniformly random hash, and think
of $h(a)$ as an infinite binary fraction
$0.b_1 b_2 b_3 \cdots$. Define $\rho(x)$ to be the position of the
first $1$ bit in this expansion: $\rho(0.001\ldots) = 3$,
$\rho(0.1\ldots) = 1$, and so on. Equivalently, $\rho(x)$ counts the
number of leading zeros and adds one.

For a uniformly distributed $h(a)$ the bit $b_k$ is an independent
fair coin, so
\[
  \Prob{\rho(h(a)) \ge k} \;=\; \Prob{b_1 = b_2 = \cdots = b_{k-1} = 0}
  \;=\; 2^{-(k-1)},
\]
and the indicator that $\rho(h(a)) \ge k$ is exactly the event of
seeing $k - 1$ leading zeros. The expected value of $\rho(h(a))$ over
a single draw is therefore $\sum_{k \ge 1} \Prob{\rho \ge k} =
\sum_{k \ge 0} 2^{-k} = 2$, but the tail is heavy: the probability
of seeing at least $k$ leading zeros is $2^{-k}$, so the maximum of
$\rho(h(a_1)), \ldots, \rho(h(a_n))$ across $n$ independent items
grows as $\log_2 n$. Concretely, in $n$ independent draws of a
$\mathrm{Geometric}(1/2)$ variable, the maximum is approximately
$\log_2 n$ with standard deviation $O(1)$: each unit increase in the
threshold halves the probability of exceeding it, so the maximum
sits at the value where the expected count of exceedances drops from
above one to below one.

The trick of HyperLogLog is to invert this relationship. Track only
the running maximum of $\rho(h(a))$ over the items seen, and read
off an estimate of $\log_2 n$ from it. The data the algorithm needs
to maintain is a single small integer (the current maximum $\rho$),
not the items themselves.

\subsection*{A single-register estimator}

The crudest version of the algorithm is the obvious one. Maintain a
single register $R$, initialized to zero. On each stream item $a$,
compute $\rho(h(a))$ and set $R \gets \max(R, \rho(h(a)))$. At query
time, report
\[
  \hat n \;=\; 2^R.
\]

This single-register estimator captures the right scaling: as the
true cardinality $n$ grows, $R$ grows as $\log_2 n$, and $2^R$ grows
as $n$. But the multiplicative variance is enormous. The maximum of
$n$ independent geometric variables has standard deviation on the
order of one bit, so the estimate $2^R$ has multiplicative standard
deviation on the order of $2^{O(1)}$, that is, a constant factor of
roughly two in either direction. An estimate that is off by a factor
of two is useful for an order-of-magnitude check and useless for
anything quantitative. The fix is averaging: run many independent
copies of the algorithm and combine them in a way that pulls the
estimate toward its mean.

\subsection*{Multiple registers (the construction)}

Running $m$ truly independent hash functions and tracking $m$
independent registers would solve the variance problem at $m$-fold
cost in hash evaluations. HyperLogLog avoids that overhead by reusing
a single hash via \emph{stochastic averaging}: each item is routed to
one of $m$ buckets by the first $\log_2 m$ bits of its hash, and only
the remaining bits feed the leading-zero count for that bucket. The
$m$ registers then play the role of independent observations of the
single-register estimator, each on a stream of expected length $n/m$,
without ever evaluating a second hash.

\begin{definition}[HyperLogLog sketch]
\label{def:hyperloglog}
Fix a register count $m = 2^p$ for some integer $p \ge 4$, and a
hash $h : \Universe \to \{0, 1\}^{32}$ (or 64; the bit width sets the
maximum representable cardinality). The \emph{HyperLogLog sketch} is
a register array $M[0..m-1]$ of small non-negative integers,
initialized to zero, together with the operations:
\begin{itemize}
  \item \textsc{Update}$(a)$: compute $h(a)$; split it into the
    first $p$ bits $j$ (the bucket index) and the remaining bits
    $w$; set
    $M[j] \gets \max(M[j], \rho(w))$.
  \item \textsc{Estimate}$()$: return
    \[
      \hat n \;=\; \alpha_m \cdot \frac{m^2}{\sum_{j=0}^{m-1} 2^{-M[j]}},
    \]
    where $\alpha_m$ is the bias-correction constant given below.
\end{itemize}
The constants are $\alpha_{16} = 0.673$, $\alpha_{32} = 0.697$,
$\alpha_{64} = 0.709$, and
$\alpha_m = 0.7213 / (1 + 1.079/m)$ for $m \ge 128$.
\end{definition}

\Cref{alg:hyperloglog-update-estimate} states the two operations in
pseudocode.

\begin{algorithm}[ht]
  \caption{HyperLogLog update and estimate}
  \label{alg:hyperloglog-update-estimate}
  \SetKwInOut{Input}{input}
  \SetKwInOut{Output}{output}
  \SetKwProg{func}{procedure}{}{}
  \SetKwProg{fn}{function}{}{}

  \func{\textsc{Update}($a$)}{
    $x \gets h(a)$\;
    $j \gets$ first $p$ bits of $x$\;
    $w \gets$ remaining bits of $x$\;
    $M[j] \gets \max(M[j], \rho(w))$\;
  }
  \BlankLine
  \fn{\textsc{Estimate}()}{
    $Z \gets \sum_{j=0}^{m-1} 2^{-M[j]}$\;
    \Return $\alpha_m \cdot m^2 / Z$\;
  }
\end{algorithm}

Three pieces of the estimator deserve comment. The sum $Z = \sum_j
2^{-M[j]}$ is the \emph{harmonic mean} (up to a factor of $m$) of the
per-bucket estimates $2^{M[j]}$. Harmonic averaging suppresses the
contribution of an outlier register that, by bad luck, accumulated
an unusually long run of leading zeros and would otherwise pull the
arithmetic mean of the per-bucket estimates upward by a constant
multiplicative factor. The factor $m^2$ converts an estimate of
$n/m$ per bucket into an estimate of $n$ total, and the bias-correction
$\alpha_m$ removes a small constant-factor bias introduced by the
choice of harmonic mean and the discrete geometric distribution of
$\rho$. The asymptotic form $\alpha_\infty = 0.7213\ldots$ comes
from an integral that \textcite{flajolet-et-al-2007} compute in their
Theorem 1; the small-$m$ corrections are tabulated values that fit
the same integral evaluated at finite $m$.

The space cost is $m$ registers, each large enough to hold $\rho$ of
the bit-string after the prefix. For a 32-bit hash and $p = 8$
($m = 256$), each $M[j]$ takes values in $\{0, 1, \ldots, 24\}$ and
fits in 5 bits. The whole sketch is $256 \cdot 5 = 1280$ bits, or
160 bytes, regardless of how many items the stream contains.

\subsection*{The error bound}

The bias correction is calibrated to make $\mathbb{E}[\hat n] = n$
to first order in $1/m$. The fluctuations of $\hat n$ around $n$ are
governed by the central limit theorem applied to the registers.

\begin{proposition}[HyperLogLog standard error]
\label{prop:hyperloglog-error}
For a HyperLogLog sketch with $m \ge 128$ registers built from a
uniformly random hash, the relative standard error of the
cardinality estimate $\hat n$ satisfies
\[
  \frac{\sigma(\hat n)}{n}
  \;\approx\; \frac{1.04}{\sqrt{m}}
\]
in the regime where $n$ is large compared to $m$. With $m = 2048$
(an 11-bit bucket prefix), the relative standard error is
approximately $2.3\%$.
\end{proposition}

The constant $1.04$ is the asymptotic value of the per-register
standard-deviation factor computed by \textcite{flajolet-et-al-2007}
in their main analysis; the proof is an exact generating-function
calculation that we will not reproduce here. The dependence on
$1/\sqrt{m}$ matches the central-limit intuition: $m$ buckets give
$m$ approximately independent observations of $n/m$, and the relative
error of their averaged estimator scales as $1/\sqrt{m}$ regardless
of $n$. Doubling the relative accuracy quadruples the register count;
this is the space-accuracy duality of the sketch.

Two regime boundaries are worth noting. When $n$ is small relative
to $m$ (so that many registers are still zero), the estimator
under-counts and the algorithm switches to a linear-counting
correction: $\hat n = m \ln(m / V)$, where $V$ is the count of
zero registers. When $n$ approaches the maximum representable
cardinality (set by the hash bit width), a saturation correction
applies. Both corrections are routine and are documented in the
original paper. The bulk of the operating range is governed by
\Cref{prop:hyperloglog-error}.

\subsection*{The $\Bernoulli{T}$ type}

HyperLogLog instantiates the framework of
\Cref{ch:approximate-functions} with output type $T = \mathbb{N}$,
the same as CountMin, but with a fundamentally different channel
structure. The query is a function from streams to non-negative
integers, and its error behaviour is summarized by a channel matrix
$Q$ in the sense of \Cref{def:n-ary-channel}: rows are indexed by
the true cardinality $n$ and columns by the reported estimate
$\hat n$, with entry
\[
  Q_{n, \hat n} \;=\;
  \Prob{\textsc{Estimate}() = \hat n \Given \text{true cardinality is } n}.
\]
Unlike CountMin, the matrix is \emph{not} triangular. HyperLogLog
errs in both directions: an unlucky sequence of hash bits can produce
a register that overestimates $\log_2 n$ in some buckets while
others underestimate, and the harmonic-mean combination can shake
out either above or below $n$. The shape of each row is approximately
log-normal in $\hat n / n$: by \Cref{prop:hyperloglog-error}, the
multiplicative error has constant relative standard deviation
$1.04/\sqrt{m}$ across the operating range, and the asymptotic
analysis of \textcite{flajolet-et-al-2007} shows the limiting
distribution of $\log \hat n - \log n$ is approximately normal
with that standard deviation. Each row of $Q$ therefore concentrates
on a roughly geometric interval around $n$, symmetric on a
multiplicative scale rather than an additive one. The contrast with
CountMin's one-sided additive error is the main qualitative
distinction between the two channel matrices.

\subsection*{Merging two HyperLogLog sketches}

Two HyperLogLog sketches $M^{(1)}$ and $M^{(2)}$ built on the same
hash function and the same register count $m$ on two streams $S_1$
and $S_2$ admit a particularly clean merge:
\[
  (M^{(1)} \sqcup M^{(2)})[j]
  \;=\; \max\bigl(M^{(1)}[j], M^{(2)}[j]\bigr)
  \quad \text{for all } j \in \{0, 1, \ldots, m-1\}.
\]
The merged register array is exactly the array that would have
resulted from running \textsc{Update} on the union stream $S_1 \cup
S_2$ (treating each stream as a multiset and taking the union of
their distinct items). The argument is one line: for each bucket
$j$, the bucket-$j$ register on the union stream is the maximum of
$\rho(w)$ over all items in $S_1 \cup S_2$ whose hash prefix is $j$,
which is the maximum of the per-stream bucket-$j$ maxima. The merge
is idempotent in duplicate items: an item appearing in both streams
contributes the same $\rho(w)$ to its bucket in both sketches, and
$\max$ of two equal values is that value.

In the language of \Cref{ch:approximate-functions}, this merge is a
$\Bernoulli{\mathbb{N}}$ composition. The channel matrix of the
merged sketch is the channel matrix that
\Cref{thm:n-ary-kronecker} would assign to the joint query of the
two stream sketches under the cardinality-of-union operation on
multisets. As with CountMin, sharing the hash function across the
two sketches is the discipline that lets independently constructed
sketches add up to a correct sketch of the union: were the two
sketches built on independent hashes, the merge would be a coarse
upper bound on the union cardinality with no clean error analysis.

\subsection*{A worked example}

Take $m = 256$ (so $p = 8$) and simulate a stream of $n = 10000$
distinct items, each hashed with a 32-bit hash. The expected relative
standard error from \Cref{prop:hyperloglog-error} is
$1.04/\sqrt{256} = 0.065$, so a single run typically returns an
estimate within $\pm 650$ of $10000$.

After processing the stream, each register $M[j]$ records the
maximum $\rho$ of the $\approx 10000 / 256 \approx 39$ items that
landed in bucket $j$. The expected maximum of $39$ independent draws
of $\rho$ is approximately $\log_2 39 \approx 5.3$, so most of the
registers contain values in $\{4, 5, 6, 7\}$, with a small minority
at $3$ or $8$ from the tails of the maximum-of-geometrics
distribution. A representative register array (verified in the
chapter notebook with seed 0) is concentrated in this range, with
distribution
\[
  M \;\approx\; (5, 5, 6, 4, 7, 5, 6, 5, 4, 6, 5, 7, 5, 6, 5, 4, \ldots)
\]
for the first sixteen entries, and a full sum $\sum_j 2^{-M[j]}$ on
the order of $256 / 39 \approx 6.5$. The estimate is
\[
  \hat n \;=\; \alpha_{256} \cdot \frac{256^2}{Z}
  \;\approx\; 0.7182 \cdot \frac{65536}{6.5}
  \;\approx\; 7240,
\]
which is in the wrong tail by about $2.8$ standard deviations on a
single seed; averaging over a hundred independent seeds returns an
empirical mean estimate of $9970$ and an empirical relative standard
deviation of $0.063$, both within the expected range from
\Cref{prop:hyperloglog-error}. The example reinforces the point that
the headline bound is a standard error, not a deterministic bound:
any single sketch can land a few standard deviations away from $n$,
and the value of the sketch is that the spread shrinks as
$1/\sqrt{m}$ with the register count.

\subsection*{Looking ahead}

HyperLogLog estimates the cardinality of a single multiset.
\Cref{sec:minhash-lsh} turns to a similarity question on two sets:
estimate the Jaccard similarity $J(A, B) = |A \cap B|/|A \cup B|$
from compact signatures of $A$ and $B$, without recomputing either
intersection or union explicitly. The construction is again hash-based,
again maintains a small array per set, and again admits a clean
$\Bernoulli{T}$ analysis, this time with $T = [0, 1]$ and a symmetric
two-sided error bound from Hoeffding's inequality.
